\begin{figure}[htbp]
    \centering
    \begin{tikzpicture}[>=Stealth, thick, scale=0.7]

        % --- CONFIGURATION ---
        \def\R{6}          % Radius of the sphere (length of vp)
        \def\angTheta{45}  % Angle theta
        \def\angAlpha{65}  % Angle of vx relative to vertical
        \def\lenVx{4.5}    % Length of vx vector

        % Coordinates
        \coordinate (Origin) at (0,0);           % The "North Pole" (Tangent point)
        \coordinate (Center) at (0, -\R);        % The Center of the sphere

        % --- AXES ---
        % Vertical dashed axis (upwards)
        \draw[dashed, thick] (0, 0) -- (0, 2.5);

        % Horizontal dashed axis (Tangent Space)
        \draw[dashed, thick] (-2, 0) -- (7, 0) node[right] {tangent space};


        % --- TOP PART: PROJECTION ---

        % Vector vx
        \coordinate (Vx) at (\angAlpha:\lenVx); % Defined by angle alpha
        % Note: In the image, vx is (x,y). Let's approximate the slope visually.
        \coordinate (Vx_End) at (5, 2);
        \draw[->, ultra thick] (Origin) -- (Vx_End) node[midway, above left] {$\mathrm{v_x}$};

        % Vector vt (Projection on horizontal)
        \coordinate (Vt_End) at (5, 0);
        \draw[->, ultra thick] (Origin) -- (Vt_End) node[midway, above] {$\mathrm{v_t}$};

        % Projection line 'a'
        \draw[dashed, thick] (Vx_End) -- (Vt_End) node[midway, right] {$\mathrm{a}$};

        % Angle Alpha
        \coordinate (Y_Axis_Point) at (0, 2);
        \pic [draw, pic text=$\alpha$, angle radius=0.8cm, angle eccentricity=1.3] {angle = Vx_End--Origin--Y_Axis_Point};


        % --- BOTTOM PART: GEODESIC UPDATE ---

        % Calculate the "Child" point on the sphere surface
        % x = R * sin(theta)
        % y = -R + R * cos(theta) (relative to origin)
        % But we draw from Center (0, -R).
        % Target point S relative to Center is (R sin theta, R cos theta).

        \coordinate (S) at ({\R*sin(\angTheta)}, {-\R + \R*cos(\angTheta)});
        \coordinate (M) at (0, {-\R + \R*cos(\angTheta)}); % The corner of the triangle

        % 1. Vertical Component (vp * cos theta)
        % Draws from Center to M
        \draw[->, ultra thick] (Center) -- (M) node[midway, left=5pt] {$\mathrm{v_p} \cdot \cos\theta$};

        % 2. Horizontal Component (vt/|vt| * sin theta)
        % Draws from M to S
        \draw[->, ultra thick] (M) -- (S) node[midway, above] {$\frac{\mathrm{v_t}}{\|\mathrm{v_t}\|} \sin\theta$};

        % 3. Result Vector (vc)
        % Draws from Center to S
        \draw[->, ultra thick] (Center) -- (S) node[midway, below right] {$\mathrm{v_c}$};

        % 4. Parent Vector (vp) label logic
        % The line goes from Center to Origin.
        % We need an arrow head at the Origin.
        \draw[->, ultra thick] (Center) -- (Origin); % completing the vertical line
        \node[above=40pt, left=5pt] at (0, -2.5) {$\mathrm{v_p}$}; % Label roughly in the middle of the top section

        % 5. Angle Theta
        \pic [draw, pic text=$\theta$, angle radius=1.0cm, angle eccentricity=1.3] {angle = S--Center--M};

        % 6. The Geodesic Arc
        % Arcs from Origin (top) down to S.
        % To draw an arc in TikZ, we specify start angle, end angle, and radius.
        % Start: 90 degrees (relative to Center). End: 90 - theta.
        \draw[thick] (0,0) arc (90:{90-\angTheta}:\R);
    \end{tikzpicture}
    \caption{Geometric diagram of the brownian motion simulation algorithm. }
    \label{fig:algorithm-bm}
\end{figure}